%!TEX root = ../thesis.tex
%*******************************************************************************
%****************************** Background Chapter *********************************
%*******************************************************************************

\chapter{Background}
\ifpdf
    \graphicspath{{Chapters/Background/Figs/}{Chapters/Background/Figs/}{Chapters/Background/Figs/}}
\else
    \graphicspath{{Chapters/Background/Figs/}{Chapters/Background/Figs/}}
\fi
To build the foundation of our approach, we present the usage of crowdsourcing in software development (see section \ref{background:section:crowdsourcing}), UI Prototyping (see section \ref{background:section:uiprototyping}), Low and No code (see section \ref{background:section:lowcode}), Model Based Software Engineering (MBSE) (see section \ref{background:section:mbse}), Task Based Usability Testing (see section \ref{background:section:task}), and Experimental Product Design and define Design Principles (DP) (see section \ref{background:section:experimentproduct}).

%********************************** % Crowdsourcing of Software Products **************************************
\section{Crowdsourcing of Software Products}
\label{background:section:crowdsourcing}
Iterative feedback from potential customers can help improve the development of software products \cite{article:lean:eric}.
To do that, we can use crowdsourcing.
Crowdsourcing refers to outsourcing value-creating activities from a company by an open call to a large, undefined group of users to get feedback \cite{article:crowdsourcing:leimeister}.
The word crowdsourcing is a combination of crowd and outsourcing.
Crowdsourcing often involves less specialized and more generalized groups of participants than outsourcing \cite{article:crowdsourcing:estelles}.
Some advantages of crowdsourcing include lowered costs, improved speed, quality, flexibility, and scalability \cite{article:crowdsourcing:prpic}.
Researchers have used crowdsourcing in many research approaches, including \textit{crowd testing, crowd funding, crowd ideation, crowd logistic, crowd production, crowd promotion,} and \textit{crowd support} over the last few years \cite{article:crowdsourcing:durward}.
In our approach to finding a solution, we would focus more on crowd-testing and crowd-ideation.

\paragraph{Crowd Testing:}
The companies use crowd-testing to evaluate different running software products with the users.
A growing trend in software testing is crowd testing, which utilizes the benefits, effectiveness, and efficiency of crowdsourcing and cloud platforms \cite{article:crowdsourcing:latoza}.
Crowd testing is considered when the software is more user-centric: i.e., software with a broad user base whose success is evaluated by user input.
CrowdStudy \cite{article:crowdsourcing:nebeling} is a method that enables developers to assess the usability of their web interfaces using crowd workers from Amazon Mechanical Turk\footnote{Amazon Mechanical Turk: \url{https://www.mturk.com}}.
CrowdCrit \cite{article:crowdsourcing:luther} is another tool that uses Amazon Mechanical Turk to support designers in validating created posters in the form of uploaded images.
Similarly, \textit{Interactive event-flow graphs} and \textit{GUI-level (Graphical User Interface) guidance} \cite{article:crowdsourcing:chen} are the two techniques to increase crowd testers' coverage for GUI using crowd-testing.

\paragraph{Crowd Ideation:}
Design can be infused with creativity by online crowds, but using traditional strategies to harness them, such as large-scale ideation platforms, requires organization and time \cite{article:crowdsourcing:andolina}.
Hence, crowd ideation is used to build new and improved versions of existing software product ideas with the consumers.
Under manipulations of task complexity, idea representation, and procedural guidance, Shixuan Fu et al. \cite{article:crowdsourcing:fu} examine how cognitive load is altered during idea generation and convergence with crowds.
ERICA \cite{article:crowdsourcing:erica} is a tool that uses expert knowledge to validate diverse crowd answers.
Crowdboard \cite{article:crowdsourcing:andolina} is a tool used to engage crowds in real-time brainstorming, concept mapping, and other design processes at an early stage of the design process.
There were, however, no approaches that directly addressed prototype application areas.

%********************************** % UI Prototyping **************************************
\section{UI Prototyping}
\label{background:section:uiprototyping}
User Interface prototyping is an evaluation and testing technique according to User-Centred Design (UCD) methodology since the 1990s \cite{article:prototyping:preece}. 
The evaluation of prototypes by users is a fundamental part of all iterative approaches for IT project management, especially agile methodologies \cite{article:prototyping:schwaber}.
Therefore, designing UI prototypes enables designers and stakeholders to communicate more effectively.
An interactive prototype helps visualize design concepts and communicate new requirements and expectations about a prospective system.
Thus, usability testing to validate user requirements and prototype functionality is part of the evaluation process for UI prototypes.

\begin{figure}[htbp!]
\centering    
\includegraphics[width=1.0\textwidth]{minion.png}
\caption[Minion]{This is just a long figure caption for the minion in Despicable Me from Pixar}
\label{fig:minion}
\end{figure}

%********************************** % Low code **************************************
\section{Low Code / No Code}
\label{background:section:lowcode}

%********************************** % Model-Based Software Engineering **************************************
\section{Model-based Software Engineering}
\label{background:section:mbse}

\begin{enumerate}
\item The first topic is dull
\item The second topic is duller
\begin{enumerate}
\item The first subtopic is silly
\item The second subtopic is stupid
\end{enumerate}
\item The third topic is the dullest
\end{enumerate}

%********************************** % Task based Usability Testing **************************************
\section{Task-based Usability Testing}
\label{background:section:task}

%********************************** % Experimental Product Design **************************************
\section{Experimental Product Design}
\label{background:section:experimentproduct}

\begin{itemize}
\item The first topic is dull
\item The second topic is duller
\begin{itemize}
\item The first subtopic is silly
\item The second subtopic is stupid
\end{itemize}
\item The third topic is the dullest
\end{itemize}

\section*{Description}
\begin{description}
\item[The first topic] is dull
\item[The second topic] is duller
\begin{description}
\item[The first subtopic] is silly
\item[The second subtopic] is stupid
\end{description}
\item[The third topic] is the dullest
\end{description}

% \clearpage

\begin{landscape}

\section{Landscape}
I can cite Wall-E (see Fig.~\ref{fig:WallE}) and Minions in despicable me (Fig.~\ref{fig:Minnion}) or I can cite the whole figure as Fig.~\ref{fig:animations}


\begin{figure}
  \centering
  \begin{subfigure}[b]{0.3\textwidth}
    \includegraphics[width=\textwidth]{minion.png}
    \caption{Tom and Jerry}
    \label{fig:TomJerry}   
  \end{subfigure}             
  \begin{subfigure}[b]{0.3\textwidth}
    \includegraphics[width=\textwidth]{minion.png}
    \caption{Wall-E}
    \label{fig:WallE}
  \end{subfigure}             
  \begin{subfigure}[b]{0.3\textwidth}
    \includegraphics[width=\textwidth]{minion}
    \caption{Minions}
    \label{fig:Minnion}
  \end{subfigure}
  \caption{Best Animations}
  \label{fig:animations}
\end{figure}

\end{landscape}
